\section{Deep Predictive Learning}

\subsection{Concept of Predictive Coding}
The machine learning framework involves learning an optimal model with a large amount of data and ensuring that the model works appropriately on new data, i.e., has maximum generality. This concept is the same for object recognition and reinforcement learning. However, obtaining an optimal model is challenging, considering the impracticality of predicting which state will give maximum generalizability. Especially robots can be assumed to be constantly exposed to unexpected situations to operate in the real world. Therefore, adopting an approach that considers learning as imperfect and deals with this error is necessary. A possible framework is Friston's free energy principle~\cite{friston2012perception}, given by 

\begin{equation}
    FreeEnergy = -log(x|z) + D_{KL}[q(z|x)||p(z)] \label{eq_free}
\end{equation}

This equation corresponds to the negative evidence lower bound (ELBO) in general machine learning models, which is not new. However, based on the ``interpretation'' of this equation and its multiple parameters, we can consider designing an embodiment system.

The brain is a predictor and has as a prior belief $p(z)$ about $z$ that causes the sensation $x$. Here, perception $z$ includes the hidden states (causes) of exteroceptive sensations such as vision and touch and the hidden states of proprioceptive sensations (e.g., joints) and endoceptive sensations related to motor and visceral sensations. The brain predicts sensation $x$ based on perception $z$ and receives the current proprioceptive sensory input $x$ to obtain the posterior belief $q(z|x)$. A new prediction for the next sensation $x$ is generated based on the updated posterior belief. The prediction generation and belief updating are performed continuously. The predicted sensation $x$ includes a prediction of accuracy (the inverse of the variance of the sensation $vx$). The first term in equation \ref{eq_free} refers to the prediction error (Accuracy) weighted by this accuracy, and the second term refers to the effect of the complexity of the prior belief (model) structure (Complexity).

Strategies to reduce prediction error are called ``predictive coding,'' and four main strategies are possible.

First, learn by training the entire model that generates the sensory predictions and changing the internal parameter $w$ according to the error. However, this approach is costly and challenging to implement in time-constrained situations. Furthermore, there may be cases where frequent model modifications are inappropriate.

Second, modify the posterior belief $q(z)$. For example, when humans perceive the environment, they may comfortably interpret the real world according to the prior images in their brains and are unaware of the error. This method uses predictive biases about the world to adapt to changes in the real world by creating a virtual perception that differs from reality.

Third, use the body to influence the external real world by changing sensory $x$ to reduce prediction error. We achieve this by manipulating the environment or changing one's perspectives. This strategy is essential for autonomous agents, such as robots, with bodies in the real world.

Fourth, reduce the apparent error by setting a smaller prediction accuracy of the sense $x$, $1/vx$, and multiplying it by the ``weight'' of the prediction error.
This is related to the concept of ``attention'' discussed below.




\subsection{Comparison with Reinforcement Learning}
This section discusses the main differences between elementary reinforcement learning and predictive coding. 

One of the distinctive characteristics of reinforcement learning is that the learning goal can be set freely and abstractly (e.g., scalar values).
In contrast, the main goal in predictive coding is to minimize the prediction error in the model's connection between the body and the world.
Therefore, in principle, it is difficult to set other learning goals in predictive coding unless special efforts are made.

In model-based reinforcement learning, the ``world model'' is learned through repeated interactions between actions and states.
In \cite{ha2018world}, actions and states are learned as a set of abstracted latent spaces; the idea is similar to predictive coding in that the world is the target of the prediction.
However, the world model in reinforcement learning searches for the action policy that can achieve the maximum reward in the model by trial and error after sufficient learning has progressed.
Thus, the optimal behavior based on the learned world model can be acquired in reinforcement learning.

In contrast, the predictive coding model assumes an open world and that learning the model of the environment is continuous.
With insufficient learning, the model generates actions that are expected to have minor future prediction errors (expected free energy).
Then, by repeating the experience of reaching the target state, the sequence of actions becomes ``habituated'' to the prior belief $p(x)$.
The generated series of behavioral sequences (virtual expectations) is based on trial and error in this habituation process.
However, the learned motions need not be optimal for any given objective measure.

For example, right-handed people pick objects in front of them using their right hand, even if the object is slightly to the left.
They do not force themselves to pick it up with their left hand. This behavior is habitual, i.e., they choose an action with high predictive accuracy.
In the next section, we discuss implementing a mechanism that partially reflects this idea of predictive coding in a robot.


\subsection{Physicality Perspective}
An interesting consideration can be drawn from discussing physicality as the norm, as exemplified by behavior-based robotics~\cite{brooks1991intelligence}.
Here, it is assumed that intelligent physical behavior emerges from the coupling (mechanical attraction) of the ``environment-body-nervous system.''
In particular, it emphasizes the interaction between the body and the environment.
For example, in the passive walking mechanism, ``walking'' emerges from the interaction between a cleverly designed bipedal (or multipedal) articulation model and a slope.
We also focus on the various ``functions'' that emerge from the interaction between the skillfully designed body and the environment, including the body's flexibility, mechanism, and sensor placement.
In other words, the body itself implicitly retains the motion and actions.
At this point, the nervous system is considered to perform only minimal functions, in an example of passive walking, the oscillator system of the body is enough to execute the action.


\subsection{Dynamic Cognitive System Perspective}\label{sec_cognitive}
Given the ``environment-body-nervous system'', the nervous system can have multiple attractor states.
An example is the discussion of dynamic cognitive systems based on the dynamical systems perspective proposed by Tani ~\cite{tani1996model}.
Here, the brain is assumed to be an input-output and predictive system for the world (Figure \ref{fig_attractor}).
In interaction with the body, the brain sometimes forms a specific attractor structure with the world, and when the discrepancy grows, the connection is broken.
The brain then transitions to a new attractor, considered a phase transition, sometimes through a chaotic state.
As this process repeats, cognition never converges but transitions dynamically.
Tani et al. realized this process using a recurrent neural network (RNN) and a real robot system and have been continuously doing exciting research to understand the cognitive processes of the brain.
Although this study was before the introduction of predictive coding and the free energy principle, the concepts have much in common.


%-------------------- Fig --------------------
\setlength\textfloatsep{5pt}
\begin{figure}[t]
\centering
\includegraphics[width=\columnwidth]{attractor.eps}
\caption{Framework of a dynamic cognitive system}
\label{fig_attractor}
\hspace{10mm}
\end{figure}    


\subsection{Examples of Deep Predictive Learning} \label{sec_attention}
With the above ideological background, we believed that the ``neural system'' part of the ``environment-body-neural system''coupling could be configured in a deep learning model to couple the robot with the world to perform higher dimensional and more complex tasks. The first model was published in 2014 \cite{noda2014multimodal}, and since then, this area of research has continued to expand.

Deep predictive learning combines multiple deep learning models to build a sensory and motion prediction system. For example, convolutional autoencoder and RNNs are combined to build a sensory-motor prediction system coupled with the environment and body. The following are some critical structural points of the deep predictive learning model. The sensation and motion are integrated into a single predictive system. In a typical engineering application, mapping is learned with sensation as the input and motion as the output. However, in the deep predictive learning framework, sensation and motion form a single perceptual system that is both the input and the object of prediction. The framework is based on learning this perceptual continuum as an experience. Using RNNs with different time scales is also essential, as shown in figure \ref{fig_mtrnn}. A hierarchical approach is used where the lower layer, closer to the senses, predicts the early (short time constant) dynamics, while the upper layer predicts the late (long time constant) dynamics. Methods such as multiple-time scale RNN (MTRNN)\cite{yamashita2008emergence} and long short-term memory (LSTM)\cite{Hochreiter1997nc} are used to implement this, and the interactions help minimize the overall prediction error. Specifically, the prediction error of the senses is absorbed by the state changes in the lower layers. Here, the upper layer generates predictions based on long-term expectations (intentions) and provides constraints to the lower layer to maintain this state. However, if the lower layer cannot absorb the prediction errors of the senses, the errors affect the upper layer, changing the state of the upper layer (change in intention). Thus, the structure in which the upper and lower layers interact allows generative adaptive behavior while maintaining consistency. From a hierarchical perspective, the predictive coding inspired RNN (PV-RNN) by Tani et al. is the latest predictive coding model that allows updating the prior belief $p(z)$ at each level, including the upper and lower layers)\cite{ahmadi2019novel}. This model deals aggressively with prediction error minimization using real-time error backpropagation (called regression) and has high theoretical consistency. Although some problems, such as computational cost, exist, applying this model to real robots with multidimensional degrees of freedom can be beneficial. In addition, a recent study investigated the equivalence of hierarchical neural networks with recursive coupling and predictive coding\cite{isomura2022canonical}.




%-------------------- Fig --------------------
\setlength\textfloatsep{5pt}
\begin{figure}[t]
\centering
\includegraphics[width=60mm]{S-MTRNN.eps}
\caption{Hierarchical structure of recursively coupled predictors with different time scales}
\label{fig_mtrnn}
\hspace{10mm}
\end{figure}


\newpage

\setlength\textfloatsep{5pt}
\begin{figure}[t]
\centering
\includegraphics[width=\columnwidth]{attention_bag.eps}
\caption{Bag opening task using an attention mechanism}
\label{fig_attention_bag}
\hspace{10mm}
\end{figure}   


A more detailed structure and implementation methods of deep predictive learning are described in Section \ref{sec_implementation}. The following gives an overview of deep predictive learning using a concrete example. Figure \ref{fig_attention_bag} shows a robot~\footnote{https://youtu.be/SC0BfTlnYmw} that can open the zipper of a cloth bag, which the authors presented in ~\cite{ichiwara2022contact, ito2022integrated}. Here, the cloth bag, the robot arm (with a camera and tactile sensors), and the deep learning model are retracted so that the perceptual prediction (posterior belief) $q(z)$ in the model is modified in real-time (i.e., the predicted image, tactile state, and motion are modified). The input sensation $x$ is also modified by motion modification, and the real-time model prediction error is reduced. In conventional control technology, the trajectory of the robot arm is given in advance as a target for real-time error correction. However, in the ``zipper opening task'' discussed in this paper, generating an appropriate trajectory in real-time according to the state of the bag and robot hand is necessary. Furthermore, as this task starts with grasping a randomly placed zipper, the grasping position and posture constantly change, increasing the task's difficulty level. In this study, direct teaching is used to train the robot to open the zipper multiple times to give the robot a real experience with the task we want to convert into a habit. The number of times the movement is taught is 36, which is extremely small compared to common learning methods.


\begin{description}
\item[Trial and error]\mbox{}\\
Owing to online motion generation by the real robot using the trained model, granular trajectory modification and motion switching were based on a trial-and-error process.
Specifically, this means that the robot generates a motion based on the prediction (prior belief) $p(x)$. The motion is then modified in real-time (i.e., a transition to another motion attractor) using a retraction based on the prediction error.
The trial and error is not random but a modification of the predicted motion based on past experiences.


\item[Bottom-up and top-down attention]\mbox{}\\
By introducing an ``attention mechanism,'' the model becomes robust to sudden changes in lighting conditions, placement of multiple objects as noise, and movement of the target object, among others. Specifically, based on verbal attention instructions, we confirm the phenomenon of ``disappearance of backgrounds other than the attended object'' in the generated predictive images (posterior beliefs) ~\cite{ito2022integrated}. It is important to note that attention here is created by ``prediction.'' Normally, attention in image processing is generated from the ``current input image'' in a bottom-up manner. It involves selecting features such as luminance, color, and orientation using saliency-based feature selection of features or by generating a mask, key, query, and value from the image to maximize the detection reward. The specific implementation of top-down attention prediction is under debate. In the zipper-opening task, several points of attention are predicted in the image. In other words, we limit the image regions that can be attended to and ignore the rest. The loss function is a prediction error that includes all visual, tactile, and motor components; thus, the points of attention are self-organized into predictable parts that change (can be manipulated) depending on one's actions.
\end{description}
